\chapter{Martin Hairer}\label{ch2014martinhairer:martin-hairer-fields-medal-2014-a-comprehensive-report-on-his-work-and-impact}

\begin{center}\rule{0.5\linewidth}{0.5pt}\end{center}

\textbf{Prepared:} July 2026 \textbf{Subject:} Sir Martin Hairer KBE FRS (b. November 14, 1975, Geneva, Switzerland) \textbf{Position:} Professor of Mathematics, EPFL and Imperial College London \textbf{Doctoral Advisor:} Jean-Pierre Eckmann (University of Geneva) \textbf{Fields Medal Citation:} \emph{``For outstanding contributions to the theory of stochastic partial differential equations, and in particular for the creation of a theory of regularity structures for such equations.''}

\begin{center}\rule{0.5\linewidth}{0.5pt}\end{center}

\begin{prereqs} Probability theory (Brownian motion, stochastic processes), PDE theory (heat equation, Sobolev spaces), basic functional analysis, rough path theory. \end{prereqs}

\begin{objectives} Explain regularity structures and their role in stochastic PDEs; describe the KPZ equation and its solutions; understand the connection between rough paths and renormalization in infinite dimensions. \end{objectives}

\section{Problem History}

Stochastic partial differential equations (SPDEs) arise naturally in statistical mechanics, quantum field theory, and fluid dynamics whenever random fluctuations are modeled as space-time white noise $\xi$. A typical SPDE takes the form
\[
\partial_t u = \Delta u + F(u) + \xi,
\]
where $\xi$ is a distribution, not a function. In one spatial dimension, the linear equation has a well-defined solution (the stochastic heat equation), but in two or three dimensions, nonlinear terms like $F(u) = u^3$ or $(\partial_x h)^2$ become products of distributions---objects that have no classical meaning.

The KPZ equation, introduced by Kardar, Parisi, and Zhang in 1986,
\[
\partial_t h = \partial_x^2 h + (\partial_x h)^2 + \xi,
\]
describes the stochastic growth of interfaces and belongs to a universality class encompassing particle deposition, bacterial colony growth, flame front propagation, and fluctuations in liquid crystals. The KPZ universality class is characterized by scaling exponents (roughness $\chi = 1/2$, dynamic $z = 3/2$) independent of microscopic details. Despite its physical importance, the nonlinear term $(\partial_x h)^2$ is ill-defined when the solution $h$ is a distribution.

Physicists had long known that these equations require \emph{renormalization}---the subtraction of infinite constants---analogous to the procedure in quantum field theory. Before Hairer's work, no rigorous mathematical framework existed for this renormalization in the SPDE setting.

The only rigorous approach to the KPZ equation was through the Hopf-Cole transformation $Z = e^h$, which converts it into the linear stochastic heat equation $\partial_t Z = \partial_x^2 Z + Z\xi$. However, this trick fails for coupled KPZ systems, KPZ with multiplicative noise, and the dynamical $\Phi^4_3$ model $\partial_t\Phi = \Delta\Phi + C\Phi - \Phi^3 + \xi$, which describes the stochastic quantization of the $\Phi^4$ quantum field theory in three spatial dimensions.

Terry Lyons's theory of rough paths (1990s) provided a pathwise theory of differential equations driven by irregular signals in one-dimensional time. The central idea is to augment a rough path $X$ with its iterated integrals---the ``signature'' $\mathbb{X}_{s,t} = (1, \int_s^t dX_{r_1}, \int_s^t\int_s^{r_2} dX_{r_1}\otimes dX_{r_2}, \ldots)$---to make sense of integrals $\int Y\,dX$ when $X$ is not smooth. This theory could handle stochastic ODEs but not SPDEs, where noise varies in both space and time.

\section{How Was It Solved}

Hairer's solution \cite{hairer2014} proceeded in two stages. First, in his 2013 Annals of Mathematics paper ``Solving the KPZ equation,'' he gave the first direct construction of solutions to the KPZ equation without using the Hopf-Cole transformation. The key innovation was to treat $(\partial_x h)^2$ not as a product of distributions but as a rough path lift: one needs to keep track of additional information---the iterated integrals of the noise---to define the nonlinear term.

The construction involves a renormalization procedure: approximate the noise by smooth functions (regularization), solve the regularized equation, subtract a diverging constant from the nonlinear term, and take the limit as regularization is removed. Hairer proved that this limit exists and is independent of the approximation scheme.

Second, in his 2014 Inventiones Mathematicae paper (269--504), Hairer developed the full theory of \textbf{regularity structures}, which generalizes rough paths from one-dimensional time to multi-dimensional space-time. A regularity structure consists of: (1) an index set $A \subset \mathbb{R}$ of ``regularities,'' (2) a graded vector space $T = \oplus_{\alpha\in A} T_\alpha$ of ``modelled distributions,'' and (3) a structure group $G$ encoding the symmetries of the equation.

The centerpiece is the \textbf{Reconstruction Theorem}: a ``model'' (a consistent family of local approximations $\Pi_x: T\to\mathcal{S}'(\mathbb{R}^d)$ and connection maps $\Gamma_{xy}: T\to T$) can be reconstructed into a global distribution, analogous to how a smooth function is determined by its Taylor expansions. The solution to the SPDE is represented as a modelled distribution---a function taking values in $T$ satisfying a H\"{o}lder-type condition---and the equation is solved by a fixed-point argument.

The modelled distribution $\underline{u}$ assigns to each point $x$ an element of the model space $T$ that encodes the local behavior of the solution. The key condition is that for nearby points $x$ and $y$,
\[
\underline{u}(y) \approx \Gamma_{xy} \underline{u}(x) + \Pi_y((\underline{u}(x))_{<\alpha})
\]
where $\alpha$ is the regularity index and $(\underline{u}(x))_{<\alpha}$ denotes the truncation of $\underline{u}(x)$ to regularities less than $\alpha$. This condition ensures that the modelled distribution determines a unique global distribution via the reconstruction theorem. The fixed-point argument solves the SPDE by iterating a nonlinear map on the space of modelled distributions, and the contraction mapping principle ensures convergence when the noise is sufficiently regularized.

For the dynamical $\Phi^4_3$ equation, Hairer identified the relevant trees (rooted trees decorated with labels like $\Xi$ for the noise, $\mathcal{I}(\Xi)$ for the heat kernel applied to the noise, $\mathcal{I}(\Xi^3)$ for the cubed term), computed their divergences, defined the renormalized model by subtracting these divergences, and solved the fixed-point equation. The constant $C$ diverges as regularization is removed, but the renormalized equation has a well-defined solution.

The \textbf{renormalization group} within regularity structures consists of transformations of the model corresponding to adding counterterms. Hairer proved that the renormalized solutions are independent of the regularization scheme---a universality property essential for physical interpretation.

The algebraic structure underlying regularity structures is that of a \textbf{Hopf algebra} of rooted trees. The renormalization procedure corresponds to a renormalization group flow on this Hopf algebra, analogous to the Bogoliubov--Parasiuk renormalization in quantum field theory. The key insight is that the divergences in the SPDE correspond to specific trees in the Hopf algebra, and the counterterms are determined by the Hopf algebra structure. This connection between SPDEs and algebraic combinatorics was one of the most unexpected aspects of Hairer's theory.

The \textbf{modelled distribution} is the central object in the theory. It assigns to each point $x$ in space-time an element of the model space $T$, which encodes the local behavior of the solution. The modelled distribution satisfies a H\"{o}lder-type condition: the difference between the values at two nearby points is controlled by the distance between the points, weighted by the regularity index. This condition is the infinite-dimensional analogue of the H\"{o}lder condition for classical functions, and it ensures that the modelled distribution determines a unique global distribution via the reconstruction theorem.

\section{Impact of This Discovery}

The theory of regularity structures has created a new subfield of stochastic analysis, taken up by researchers worldwide. It provides the rigorous mathematical foundation for the stochastic quantization approach to Euclidean quantum field theory, pioneered by Parisi and Wu in 1981, in dimensions 2 and 3.

The theory resolves the long-standing problem of defining solutions to the $\Phi^4_3$ model, one of the most important examples in constructive quantum field theory. The stochastic quantization approach gives dynamics to the $\Phi^4$ Euclidean quantum field, which had previously been studied only in its static (equilibrium) form.

In quantitative finance, the theory of rough paths and regularity structures provides the mathematical tools for analyzing ``rough volatility'' models---which use fractional Brownian motion with small Hurst parameter to model financial volatility---that have become increasingly important in pricing and risk management.

Hairer's work on the stochastic Navier-Stokes equations with Jonathan Mattingly, proving unique ergodicity and developing the theory of hypoellipticity for infinite-dimensional systems, has had a major impact on the study of stochastic fluid dynamics. These results show that even with degenerate noise, the nonlinear dynamics can propagate randomness to all modes, resulting in a unique stationary measure.

The method of controlled rough paths, developed with Gubinelli---where a path $Y$ is ``controlled'' by $X$ if $Y_t - Y_s = Y'_s(X_t - X_s) + O(|t-s|^{2\alpha})$---has become a standard technique in stochastic analysis, applicable far beyond its original context.

The $\Phi^4_3$ model describes a scalar field $\Phi$ in three spatial dimensions with a quartic self-interaction. The Euclidean action is
\[
S(\Phi) = \int_{\mathbb{R}^3} \left(\frac{1}{2}|\nabla\Phi|^2 + \frac{m^2}{2}\Phi^2 + \frac{\lambda}{4!}\Phi^4\right) dx.
\]
The stochastic quantization approach adds a time variable and interprets the Euclidean field as the stationary distribution of a stochastic process:
\[
\partial_t \Phi = \Delta\Phi - m^2\Phi - \frac{\lambda}{3!}\Phi^3 + \sqrt{2}\xi,
\]
where $\xi$ is space-time white noise. The nonlinear term $\Phi^3$ is a product of distributions and is ill-defined without renormalization. Hairer's regularity structures provide the mathematical framework for making sense of this equation, by treating $\Phi^3$ as a renormalized object that absorbs the infinities arising from the product of distributions.

The renormalization procedure for the $\Phi^4_3$ model involves three trees: the ``tadpole'' tree $\mathcal{I}(\Xi)$, the ``setting sun'' tree $\mathcal{I}(\Xi^2)$, and the ``double setting sun'' tree $\mathcal{I}(\Xi^3)$. Each tree corresponds to a divergent integral that must be subtracted to obtain a well-defined equation. The renormalized equation has the form
\[
\partial_t \Phi = \Delta\Phi - m^2\Phi - \frac{\lambda}{3!}:\!\Phi^3\!: + \sqrt{2}\xi,
\]
where $:\!\Phi^3\!:$ denotes the Wick-ordered product, which is the renormalized version of $\Phi^3$. Hairer's theory proves that this renormalized equation has a unique solution that is independent of the regularization scheme.

\section{Modern Views and Improvements}

Around the same time as Hairer's work, Gubinelli, Imkeller, and Perkowski developed the theory of \emph{paracontrolled distributions}, an alternative approach to singular SPDEs using paraproducts from harmonic analysis. The two theories are complementary: paracontrolled distributions are simpler in some respects but less general than regularity structures.

Hairer and collaborators continue to extend the theory to multi-scale regularity structures for systems with multiple characteristic scales, non-parabolic equations (hyperbolic SPDEs), and equations with non-local operators. The application to quantum field theory is an active area: Hairer is a principal investigator in the Simons Collaboration on Probabilistic Paths to Quantum Field Theory.

New numerical methods for simulating singular SPDEs have been developed based on the algebraic structure of regularity structures. By incorporating the renormalization framework, numerical schemes can correctly capture the behavior of solutions. These methods have been applied to the $\Phi^4_3$ model and other singular SPDEs, producing approximations that converge to the correct limit as the mesh size decreases.

Recent work has explored connections between rough path theory and machine learning. The signature of a path provides natural features for learning tasks involving sequential data, and regularity structures may provide a framework for learning solutions to PDEs. The signature of a path $X$ is the tensor $\mathbb{X} = (1, \int dX, \int\int dX\otimes dX, \ldots)$, which encodes all iterated integrals of the path. This object has been used as a feature vector in machine learning models for time series analysis, and the connection to regularity structures suggests new ways to incorporate PDE constraints into learning algorithms.

The theory has also found applications in the study of wave turbulence, where Hairer's techniques for controlling nonlinear wave interactions have proved essential for deriving kinetic equations from first principles.

Current research frontiers include the development of regularity structures for SPDEs on Riemannian manifolds, the study of singular SPDEs with fractional Laplacians, and the rigorous construction of lattice gauge theories using renormalization techniques from regularity structures.

The application to quantum field theory is perhaps the most ambitious direction. The Yang--Mills equations in three dimensions, which describe the strong force in particle physics, are singular SPDEs that require renormalization. Hairer's theory provides the mathematical framework for making sense of these equations, and recent progress has established partial results on the construction of solutions.

The numerical analysis of singular SPDEs is another active area. Traditional numerical methods fail for equations like the $\Phi^4_3$ model because the solution is not smooth enough to be approximated by standard finite element methods. Hairer's regularity structures provide a framework for developing numerical schemes that correctly capture the renormalization structure, leading to convergent approximations.

The connection between regularity structures and the theory of rough paths continues to be explored. The rough path theory provides a pathwise approach to stochastic differential equations, while regularity structures provide a space-time approach to SPDEs. Understanding the relationship between these two theories---and developing a unified framework---is an active area of research.

The theory of regularity structures has also found applications in the study of interacting particle systems and in the theory of random geometry. The KPZ equation, which was one of the motivating examples for Hairer's work, describes the growth of random interfaces and has applications to the study of random trees, last passage percolation, and the distribution of longest increasing subsequences. The regularity structures framework provides a systematic way to study these models and to derive their scaling limits.

The study of the KPZ universality class has been enriched by Hairer's work. The KPZ equation belongs to a universality class that includes many models of random growth, such as the asymmetric simple exclusion process (ASEP), random tilings, and last passage percolation. The connection between these models and the KPZ equation is through the Brownian motion and the Tracy-Widom distribution, which describe the fluctuations of the growth process. Hairer's regularity structures provide a rigorous framework for studying these connections and for deriving the scaling limits of models in the KPZ universality class.

The development of new techniques for analyzing singular SPDEs continues to be an active area. Recent work has explored connections between regularity structures and the theory of regularity structures on manifolds, where the geometry of the underlying space introduces new complications. The theory of regularity structures for SPDEs on Riemannian manifolds has potential applications to quantum field theory and to the study of random geometry.

The application of regularity structures to the study of stochastic quantization of quantum fields has been a major focus of recent research. The stochastic quantization approach interprets Euclidean quantum fields as the stationary distribution of a stochastic process, which is described by a singular SPDE. Hairer's regularity structures provide the mathematical framework for making sense of these equations, and the resulting construction of quantum fields has implications for mathematical physics and for the foundations of quantum field theory.

The connection between regularity structures and the theory of renormalization in quantum field theory is an active area of research. The renormalization procedure in quantum field theory involves the subtraction of divergent quantities, which is analogous to the renormalization in Hairer's theory. The algebraic structure underlying both renormalization procedures is that of a Hopf algebra, and the connection between these algebraic structures has led to new insights into the mathematical foundations of quantum field theory.

The study of the KPZ universality class has been enriched by Hairer's work. The KPZ equation belongs to a universality class that includes many models of random growth, such as the asymmetric simple exclusion process (ASEP), random tilings, and last passage percolation. The connection between these models and the KPZ equation is through the Brownian motion and the Tracy-Widom distribution, which describe the fluctuations of the growth process. Hairer's regularity structures provide a rigorous framework for studying these connections and for deriving the scaling limits of models in the KPZ universality class.

The study of the hydrodynamic limits of interacting particle systems is an active area that draws on Hairer's techniques. The hydrodynamic limit describes the passage from a microscopic model of interacting particles to a macroscopic PDE, and the derivation of the hydrodynamic limit requires controlling the correlations between particles. The techniques developed by Hairer for singular SPDEs have been applied to the study of hydrodynamic limits, leading to new results on the convergence of interacting particle systems to their hydrodynamic limits.

The theory of regularity structures has also found applications in the study of interacting particle systems and in the theory of random geometry. The KPZ equation, which was one of the motivating examples for Hairer's work, describes the growth of random interfaces and has applications to the study of random trees, last passage percolation, and the distribution of longest increasing subsequences. The regularity structures framework provides a systematic way to study these models and to derive their scaling limits.

The study of the KPZ universality class has been enriched by Hairer's work. The KPZ equation belongs to a universality class that includes many models of random growth, such as the asymmetric simple exclusion process (ASEP), random tilings, and last passage percolation. The connection between these models and the KPZ equation is through the Brownian motion and the Tracy-Widom distribution, which describe the fluctuations of the growth process. Hairer's regularity structures provide a rigorous framework for studying these connections and for deriving the scaling limits of models in the KPZ universality class.

\begin{center}\rule{0.5\linewidth}{0.5pt}\end{center}

\section*{Summary}
\addcontentsline{toc}{section}{Summary}

Hairer created the theory of regularity structures, a general framework for solving singular stochastic PDEs that resolved the KPZ equation and the $\Phi^4_3$ model, revolutionizing stochastic analysis and its applications to physics and finance.
